% Based on evan.sty. Here comes the original license.
% This will go away as soon as I am done gutting this mess.

%%%%%%%%%%%%%%%%%%%%%%%%%%%%%%%%%%%%%%%%%%%%%%%%%%%%%%%%%%%%%%%%%%%%%%%%%%%%%%
%
% BOOST SOFTWARE LICENSE - VERSION 1.0 - 17 AUGUST 2003
%
% Copyright (c) 2025 Evan Chen [evan at evanchen.cc]
% https://web.evanchen.cc/ || github.com/vEnhance
%
% Available for download at:
% https://github.com/vEnhance/dotfiles/blob/main/texmf/tex/latex/evan/evan.sty
%
% Permission is hereby granted, free of charge, to any person or organization
% obtaining a copy of the software and accompanying documentation covered by
% this license (the "Software") to use, reproduce, display, distribute,
% execute, and transmit the Software, and to prepare derivative works of the
% Software, and to permit third-parties to whom the Software is furnished to
% do so, all subject to the following:
%
% The copyright notices in the Software and this entire statement, including
% the above license grant, this restriction and the following disclaimer,
% must be included in all copies of the Software, in whole or in part, and
% all derivative works of the Software, unless such copies or derivative
% works are solely in the form of machine-executable object code generated by
% a source language processor.
%
% THE SOFTWARE IS PROVIDED "AS IS", WITHOUT WARRANTY OF ANY KIND, EXPRESS OR
% IMPLIED, INCLUDING BUT NOT LIMITED TO THE WARRANTIES OF MERCHANTABILITY,
% FITNESS FOR A PARTICULAR PURPOSE, TITLE AND NON-INFRINGEMENT. IN NO EVENT
% SHALL THE COPYRIGHT HOLDERS OR ANYONE DISTRIBUTING THE SOFTWARE BE LIABLE
% FOR ANY DAMAGES OR OTHER LIABILITY, WHETHER IN CONTRACT, TORT OR OTHERWISE,
% ARISING FROM, OUT OF OR IN CONNECTION WITH THE SOFTWARE OR THE USE OR OTHER
% DEALINGS IN THE SOFTWARE.
%%%%%%%%%%%%%%%%%%%%%%%%%%%%%%%%%%%%%%%%%%%%%%%%%%%%%%%%%%%%%%%%%%%%%%%%%%%%%%

%%fakesection Some macros
%Small commands
\usepackage{amsmath,amssymb}
\usepackage{iftex}
  \usepackage[minimal]{yhmath}
  \usepackage{derivative}
\newcommand{\cbrt}[1]{\sqrt[3]{#1}}
\newcommand{\floor}[1]{\left\lfloor #1 \right\rfloor}
\newcommand{\ceiling}[1]{\left\lceil #1 \right\rceil}
\newcommand{\mailto}[1]{\href{mailto:#1}{\texttt{#1}}}
\newcommand{\ol}{\overline}
\newcommand{\ul}{\underline}
\newcommand{\wt}{\widetilde}
\newcommand{\wh}{\widehat}
\newcommand{\eps}{\varepsilon}
\newcommand{\vocab}[1]{\textbf{\color{blue}\sffamily #1}}
\providecommand{\alert}{\vocab}
\providecommand{\half}{\frac{1}{2}}
\newcommand{\catname}{\mathsf}
\newcommand{\hrulebar}{
  \par\hspace{\fill}\rule{0.95\linewidth}{.7pt}\hspace{\fill}
  \par\nointerlineskip \vspace{\baselineskip}
}
\providecommand{\arc}[1]{\wideparen{#1}}

%For use in author command
\newcommand{\plusemail}[1]{\\ \normalfont \texttt{\mailto{#1}}}

%More commands and math operators
\DeclareMathOperator{\cis}{cis}
\DeclareMathOperator*{\lcm}{lcm}
\DeclareMathOperator*{\argmin}{arg min}
\DeclareMathOperator*{\argmax}{arg max}

%Convenient Environments
\newenvironment{soln}{\begin{proof}[Soluzione]}{\end{proof}}
\newenvironment{parlist}{\begin{inparaenum}[(i)]}{\end{inparaenum}}
\newenvironment{gobble}{\setbox\z@\vbox\bgroup}{\egroup}

%From H113 "Introduction to Abstract Algebra" at UC Berkeley
\newcommand{\CC}{\mathbb C}
\newcommand{\FF}{\mathbb F}
\newcommand{\NN}{\mathbb N}
\newcommand{\QQ}{\mathbb Q}
\newcommand{\RR}{\mathbb R}
\newcommand{\ZZ}{\mathbb Z}
\newcommand{\charin}{\text{ char }}
\DeclareMathOperator{\sign}{sign}
\DeclareMathOperator{\Aut}{Aut}
\DeclareMathOperator{\Inn}{Inn}
\DeclareMathOperator{\Syl}{Syl}
\DeclareMathOperator{\Gal}{Gal}
\DeclareMathOperator{\GL}{GL} % General linear group
\DeclareMathOperator{\SL}{SL} % Special linear group
\DeclareMathOperator{\Vol}{Vol} % Special linear group

%From Kiran Kedlaya's "Geometry Unbound"
\newcommand{\abs}[1]{\left\lvert #1 \right\rvert}
\newcommand{\norm}[1]{\left\lVert #1 \right\rVert}
\newcommand{\dang}{\measuredangle} %% Directed angle
\newcommand{\ray}[1]{\overrightarrow{#1}}
\newcommand{\seg}[1]{\overline{#1}}

%From M275 "Topology" at SJSU
\DeclareMathOperator{\id}{id}
\newcommand{\taking}[1]{\xrightarrow{#1}}
\newcommand{\inv}{^{-1}}

%From M170 "Introduction to Graph Theory" at SJSU
\DeclareMathOperator{\diam}{diam}
\DeclareMathOperator{\ord}{ord}
\newcommand{\defeq}{\overset{\mathrm{def}}{=}}

%From the USAMO .tex files
\newcommand{\ts}{\textsuperscript}
\newcommand{\dg}{^\circ}
\newcommand{\ii}{\item}

% From Math 55 and Math 145 at Harvard
\newenvironment{subproof}[1][Proof]{%
\begin{proof}[#1] \renewcommand{\qedsymbol}{$\blacksquare$}}%
{\end{proof}}

\newcommand{\On}{\mathrm{On}} % ordinals
\DeclareMathOperator{\img}{im} % Image
\DeclareMathOperator{\Img}{Im} % Image
\DeclareMathOperator{\coker}{coker} % Cokernel
\DeclareMathOperator{\Coker}{Coker} % Cokernel
\DeclareMathOperator{\Ker}{Ker} % Kernel
\DeclareMathOperator{\rank}{rank}
\DeclareMathOperator{\Spec}{Spec} % spectrum
\DeclareMathOperator{\Tr}{Tr} % trace
\DeclareMathOperator{\pr}{pr} % projection
\DeclareMathOperator{\ext}{ext} % extension
\DeclareMathOperator{\pred}{pred} % predecessor
\DeclareMathOperator{\dom}{dom} % domain
\DeclareMathOperator{\ran}{ran} % range
\DeclareMathOperator{\Hom}{Hom} % homomorphism
\DeclareMathOperator{\Mor}{Mor} % morphisms
\DeclareMathOperator{\End}{End} % endomorphism


    \PassOptionsToPackage{svgnames,dvipsnames,table}{xcolor}
    \usepackage{xcolor}
      \usepackage[colorlinks=true,pdfusetitle]{hyperref}
      \hypersetup{urlcolor=RubineRed,linkcolor=RoyalBlue,citecolor=ForestGreen}
    \usepackage[nameinlink]{cleveref}

    \usepackage{amsthm}
    \usepackage{thmtools}
    \hypersetup{hypertexnames=false}
    \usepackage[framemethod=TikZ]{mdframed}
    \usetikzlibrary{shadows}
    % https://tex.stackexchange.com/a/292090/76888
    % https://github.com/marcodaniel/mdframed/issues/12
    \usepackage{xpatch}
    \xpatchcmd{\endmdframed}
      {\aftergroup\endmdf@trivlist\color@endgroup}
      {\endmdf@trivlist\color@endgroup\@doendpe}
      {}{}

    \mdfdefinestyle{mdbluebox}{%
      roundcorner=10pt,
      linewidth=1pt,
      skipabove=12pt,
      innerbottommargin=9pt,
      skipbelow=2pt,
      linecolor=blue,
      nobreak=true,
      backgroundcolor=TealBlue!5,
    }
    \declaretheoremstyle[
      headfont=\sffamily\bfseries\color{MidnightBlue},
      mdframed={style=mdbluebox},
      headpunct={\\[3pt]},
      postheadspace={0pt}
    ]{thmbluebox}

    \mdfdefinestyle{mdredbox}{%
      linewidth=0.5pt,
      skipabove=12pt,
      frametitleaboveskip=5pt,
      frametitlebelowskip=0pt,
      skipbelow=2pt,
      frametitlefont=\bfseries,
      innertopmargin=4pt,
      innerbottommargin=8pt,
      nobreak=true,
      backgroundcolor=Salmon!5,
      linecolor=RawSienna,
    }
    \declaretheoremstyle[
      headfont=\bfseries\color{RawSienna},
      mdframed={style=mdredbox},
      headpunct={\\[3pt]},
      postheadspace={0pt},
    ]{thmredbox}

    \mdfdefinestyle{mdgreenbox}{%
      skipabove=8pt,
      skipbelow=0pt,
      linewidth=2pt,
      rightline=false,
      leftline=true,
      topline=false,
      bottomline=false,
      linecolor=ForestGreen,
      backgroundcolor=ForestGreen!5,
    }
    \declaretheoremstyle[
      headfont=\bfseries\sffamily\color{ForestGreen!70!black},
      bodyfont=\normalfont,
      spaceabove=2pt,
      spacebelow=1pt,
      mdframed={style=mdgreenbox},
      headpunct={ --- },
    ]{thmgreenbox}

\mdfdefinestyle{mdorangebox}{%
			skipabove=8pt,
			linewidth=2pt,
			rightline=false,
			leftline=true,
			topline=false,
			bottomline=false,
			linecolor=BurntOrange,
			backgroundcolor=Orange!5,
		}
  
  \declaretheoremstyle[
			headfont=\bfseries\sffamily\color{Orange!70!black},
			bodyfont=\normalfont,
			spaceabove=2pt,
			spacebelow=1pt,
			mdframed={style=mdorangebox},
			headpunct={ --- },
		]{thmorangebox}

    \mdfdefinestyle{mdblackbox}{%
      skipabove=8pt,
      linewidth=3pt,
      rightline=false,
      leftline=true,
      topline=false,
      bottomline=false,
      linecolor=black,
      backgroundcolor=RedViolet!5!gray!5,
    }
    \declaretheoremstyle[
      headfont=\bfseries,
      bodyfont=\normalfont\small,
      spaceabove=0pt,
      spacebelow=0pt,
      mdframed={style=mdblackbox}
    ]{thmblackbox}

    \mdfdefinestyle{mdrosebox}{%
			roundcorner = 10pt,
			linewidth=1pt,
			skipabove=12pt,
			innerbottommargin=9pt,
			skipbelow=2pt,
			linecolor=red,
			nobreak=true,
			backgroundcolor=Mahogany!6,
		}
		\declaretheoremstyle[
			headfont=\sffamily\bfseries\color{BrickRed},
			mdframed={style=mdrosebox},
			headpunct={\\[3pt]},
			postheadspace={0pt}
		]{thmrosebox}

    \mdfdefinestyle{mdpurplebox}{%
      roundcorner=10pt,
      linewidth=1pt,
      skipabove=12pt,
      skipbelow=12pt,
      innertopmargin=9pt,
      innerbottommargin=9pt,
      linecolor=black,
      nobreak=true,
      backgroundcolor=Orchid!10,
      shadow=true,
      shadowsize=6pt,
      shadowcolor=black!30,
      frametitleaboveskip=8pt,
      frametitlebelowskip=8pt,
      frametitlebackgroundcolor=Violet!50!black,
      frametitlefont=\bfseries\sffamily\color{white},
      frametitlerule=true,
    }
    \newcommand{\listhack}{$\empty$\vspace{-2em}}

  %%fakesection Theorem setup
    \theoremstyle{definition}
    %Branching here: the option secthm changes theorems to be labelled by section
        \declaretheorem[style=thmbluebox,name=Teorema,numberwithin=section]{theorem}
      \declaretheorem[style=thmrosebox,name=Assioma,sibling=theorem]{axiom}
      \declaretheorem[style=thmbluebox,name=Lemma,sibling=theorem]{lemma}
      \declaretheorem[style=thmbluebox,name=Proposizione,sibling=theorem]{proposition}
      \declaretheorem[style=thmbluebox,name=Corollario,sibling=theorem]{corollary}
      \declaretheorem[style=thmbluebox,name=Assunzione,sibling=theorem]{assume}
      \declaretheorem[style=thmbluebox,name=Teorema,numbered=no]{theorem*}
      \declaretheorem[style=thmbluebox,name=Lemma,numbered=no]{lemma*}
      \declaretheorem[style=thmbluebox,name=Proposizione,numbered=no]{proposition*}
      \declaretheorem[style=thmbluebox,name=Corollario,numbered=no]{corollary*}
      \declaretheorem[style=thmbluebox,name=Assunzione,numbered=no]{assume*}

      \declaretheorem[style=thmgreenbox,name=Algorithm,sibling=theorem]{algorithm}
      \declaretheorem[style=thmgreenbox,name=Algorithm,numbered=no]{algorithm*}
      \declaretheorem[style=thmgreenbox,name=Claim,sibling=theorem]{claim}
      \declaretheorem[style=thmgreenbox,name=Claim,numbered=no]{claim*}

      \declaretheorem[style=thmredbox,name=Esempio,sibling=theorem]{example}
      \declaretheorem[style=thmredbox,name=Esempio,numbered=no]{example*}
      \declaretheorem[style=thmredbox,name=Esempio buffo,sibling=theorem]{exampleb}
			\declaretheorem[style=thmredbox,name=Esempio meno immediato,sibling=theorem]{examplem}
			\declaretheorem[style=thmredbox,name=Esempio più buffo,sibling=theorem]{examplebb}
			\declaretheorem[style=thmredbox,name=NON Esempio,sibling=theorem]{notexample}
			\declaretheorem[style=thmredbox,name=NON Esempio buffo,sibling=theorem]{notexampleb}

    % Osservazione-style theorems
      \declaretheorem[style=thmgreenbox,name=Osservazione,sibling=theorem]{remark}
      \declaretheorem[style=thmgreenbox,name=Osservazione,numbered=no]{remark*}
      \declaretheorem[style=thmgreenbox,name=Nota,sibling=theorem]{note}
      \declaretheorem[style=thmorangebox,name=Notazione,sibling=theorem]{notation}

      \declaretheoremstyle[
        headfont=\color{blue!40!black}\normalfont\bfseries,
        spaceabove=8pt,
        spacebelow=8pt,
        bodyfont=\normalfont
      ]{basehead}

      \declaretheoremstyle[
        headfont=\color{green!40!black}\normalfont\bfseries,
        spaceabove=8pt,
        spacebelow=8pt,
        bodyfont=\color{green!40!black}\normalfont
      ]{drkgrntheo}

    \declaretheorem[style=basehead,name=Answer,sibling=theorem]{answer}
    \declaretheorem[style=basehead,name=Answer,numbered=no]{answer*}
    \declaretheorem[style=basehead,name=Conjecture,sibling=theorem]{conjecture}
    \declaretheorem[style=basehead,name=Conjecture,numbered=no]{conjecture*}
    \declaretheorem[style=basehead,name=Definizione,sibling=theorem]{definition}
    \declaretheorem[style=basehead,name=Definizione,numbered=no]{definition*}
    \declaretheorem[style=basehead,name=Esercizio,sibling=theorem]{exercise}
    \declaretheorem[style=drkgrntheo,name=Esercizio,sibling=theorem]{drkgrnexe}
    \declaretheorem[style=basehead,name=Esercizio,numbered=no]{exercise*}
    \declaretheorem[style=basehead,name=Fatto,sibling=theorem]{fact}
    \declaretheorem[style=basehead,name=Fatto,numbered=no]{fact*}
    \declaretheorem[style=basehead,name=Problem,sibling=theorem]{problem}
    \declaretheorem[style=basehead,name=Problem,numbered=no]{problem*}
    \declaretheorem[style=basehead,name=Question,sibling=theorem]{ques}
    \declaretheorem[style=basehead,name=Question,numbered=no]{ques*}

      \Crefname{answer}{Answer}{Answers}
      \Crefname{assume}{Assumption}{Assumptions}
      \Crefname{claim}{Claim}{Claims}
      \Crefname{conjecture}{Conjecture}{Conjectures}
      \Crefname{exercise}{Exercise}{Exercises}
      \Crefname{fact}{Fact}{Facts}
      \Crefname{problem}{Problem}{Problems}
      \Crefname{ques}{Question}{Questions}

  %%fakesection Fancy section and chapter heads
{

      \renewcommand*{\sectionformat}%
        {\color{purple}\S\thesection\enskip}
      \renewcommand*{\subsectionformat}%
        {\color{purple}\S\thesubsection\enskip}
      \renewcommand*{\subsubsectionformat}%
        {\color{purple}\S\thesubsubsection\enskip}
      \addtokomafont{paragraph}{\color{orange!35!black}}
      \KOMAoptions{numbers=noenddot}
      %\usetocstyle{KOMAlike}
    }

  %%fakesection Loads a bunch of useful packages (but allow disabling)
    \usepackage{listings}
    \usepackage{mathrsfs}
    \usepackage{textcomp}
    \lstset{basicstyle=\ttfamily\small,
      backgroundcolor=\color{yellow!3!white},
      breakatwhitespace=true,
      breaklines=true,
      columns=fullflexible,
      commentstyle=\color{red!70!black},
      frame=shadowbox,
      frame=single,
      framerule=1pt,
      identifierstyle=\color{green!40!black},
      keywordstyle=\bfseries,
      keywordstyle=\bfseries\color{blue!80!black},
      numbers=left,
      numbersep=8pt,
      numberstyle=\scriptsize\sffamily\itshape\color{black!80},
      rulecolor=\color{blue!70!black},
      rulesepcolor=\color{blue!30!black},
      showstringspaces=false,
      stringstyle=\color{orange},
      tabsize=4,
      xleftmargin=15pt,
      xrightmargin=15pt,
      resetmargins=true,
    } % chktex 6
    \lstdefinelanguage{gitcommit}{
      alsoletter={:},
      morecomment=[l]{|},
      morekeywords={commit,Author:,Date:,chore,doc,edit,feat,fix,polish,style,tests,},
      sensitive=true,
    }
    \lstdefinelanguage{gitlog}{
      morekeywords={chore,doc,edit,feat,fix,polish,style,tests,},
      morecomment=[s]{[}{]}, % chktex 9
      sensitive=true,
    }

    \usepackage[shortlabels]{enumitem}
    \usepackage[obeyFinal,textsize=scriptsize,shadow,loadshadowlibrary]{todonotes}
    \usepackage{multirow}
    \usepackage{ellipsis} % don't use ugly unicode ellipsis
    \usepackage{epigraph}
    \renewcommand{\epigraphsize}{\scriptsize}
    \renewcommand{\epigraphwidth}{60ex}

    % Tiny optimizations:
    \usepackage{mathtools}
    \usepackage{microtype}
    \usepackage{xstring}
    \usepackage{wrapfig}
    \usepackage[cleanlook,british]{isodate} % day-first date
    \usepackage{tikz-cd}
    \usetikzlibrary{decorations.pathmorphing}
    \allowdisplaybreaks

    % a list I like for walkthrough's --- Drew-style parts
    \newlist{walk}{enumerate}{3}
    \setlist[walk]{label=\bfseries (\alph*)}
    % list item for MO style rubrics
    \newcommand{\worth}[1]{\def\hfill{\hskip 20pt plus 1fill}\dotfill%
      \IfEq{#1}{0}{\textbf{0~points}}%
      {\textbf{\color{blue}#1~point\IfEndWith{#1}{1}{}{s}}}%
      \par}
    \newcommand{\subworth}[1]{\def\hfill{\hskip 20pt plus 1fill}\dotfill%
      \IfEq{#1}{0}{{\footnotesize0~points}}%
      {\textbf{\footnotesize#1~point\IfEndWith{#1}{1}{}{s}}}%
      \par}
    \newlist{rubric}{enumerate}{2}
    \setlist[rubric,1]{label=\Roman*.}
    \setlist[rubric,2]{label=(\Roman{rubrici}.\alph*)}

  %%fakesection \maketitle configuration
    {% If KOMA exists. . .
    \addtokomafont{subtitle}{\Large}
    \setkomafont{author}{\Large\scshape}
    \setkomafont{date}{\Large\normalsize}
    }

\makeatletter 
  \providecommand{\thetitle}{\@title}
  \providecommand{\theauthor}{\@author}
  \providecommand{\thedate}{\@date}
\makeatother

  %%fakesection Page setup
      %\usepackage{scrlayer-scrpage}
	\pagestyle{headings}
%      \renewcommand{\headfont}{}
      \addtolength{\textheight}{5\baselineskip}
%      \setlength{\footskip}{0.5in}
%      \setlength{\headsep}{10pt}
%      \ihead{\footnotesize\textbf{\theauthor} --- \thedate}
%      \automark{section}
%      \chead{}
%      \ohead{\footnotesize\textbf{\thetitle}}
%      \cfoot{\pagemark}

%BijectiveSymbol
\newcommand{\widesim}[2][1.5]{
  \mathrel{\overset{#2}{\scalebox{#1}[1]{$\sim$}}}
}

\let\longleftrightarrow\relax % undefine the command to avoid a warning
\DeclareRobustCommand{\longleftrightarrow}{\leftarrow\bigjoinrel\rightarrow}
\newcommand\bigjoinrel{\mathrel{\mkern-7mu}}

%Other Stuff
\usepackage{leftidx}
\usepackage{stmaryrd}
\newcommand{\bigslant}[2]{{\raisebox{.2em}{$#1$}\left/\raisebox{-.2em}{$#2$}\right.}}

%Commutative Diagram
\usetikzlibrary{matrix}
\usepackage{stackengine}
\newsavebox{\tmp}

%Other Stuff of which I forgot the importance, but the they are important
\usepackage{tabularx}
\usepackage{faktor}
\newcommand{\m}{\middle|}
\newcommand{\twoheadlongrightarrow}{\relbar\joinrel\twoheadrightarrow}
\makeatletter
\newcommand{\myMod}[1]{\allowbreak \if@display \mkern 18mu \else \mkern 0mu\fi {\operator@font mod}\,\,#1} % 0mu was 12mu for the math mode in the orignial defintion
\newcommand{\myPMod}[1]{\allowbreak \if@display \mkern 18mu\else \mkern 8mu\fi \left({{\operator@font mod}\mkern 6mu #1}\right)} % replaced () by \left(\right) in respect to the original definiont
\makeatother 

%Z/nZ and F_p
\newcommand{\Z}[1]{\mathbb{Z}/{#1}\mathbb{Z}}
\newcommand{\F}[1]{\mathbb{F}_{#1}}
\newcommand{\Fpn}{\mathbb{F}_{p^{n}}}
\newcommand{\Fc}[2]{\mathbb{F}_{{#1}^{#2}}}

\DeclareMathDelimiter{\backslash}    % for double coset G\backslash H
   {\mathord}{symbols}{"6E}{largesymbols}{"0F}

\DeclareMathSymbol{\setminus}{\mathbin}{symbols}{"6E}
% Syntax: \colorboxed[<color model>]{<color specification>}{<math formula>}
\newcommand*{\colorboxed}{}
\def\colorboxed#1#{%
  \colorboxedAux{#1}%
}
\newcommand*{\colorboxedAux}[3]{%
  % #1: optional argument for color model
  % #2: color specification
  % #3: formula
  \begingroup
    \colorlet{cb@saved}{.}%
    \color#1{#2}%
    \boxed{%
      \color{cb@saved}%
      #3%
    }%
  \endgroup
}

\usepackage{xparse}

\ExplSyntaxOn
\NewDocumentCommand{\cycle}{ O{\;} m }
 {
  (
  \alec_cycle:nn { #1 } { #2 }
  )
 }

\seq_new:N \l_alec_cycle_seq
\cs_new_protected:Npn \alec_cycle:nn #1 #2
 {
  \seq_set_split:Nnn \l_alec_cycle_seq { , } { #2 }
  \seq_use:Nn \l_alec_cycle_seq { #1 }
 }
\ExplSyntaxOff

\newcommand{\pdiv}{\mid\!\mid}
\newcommand{\ps}{\mathcal{P}}
\newcommand{\rel}{\mathcal{R}}


\makeatletter
\newcommand{\skipitems}[1]{%
  \addtocounter{\@enumctr}{#1}%
}
\makeatother

\newcommand{\Cl}{\mathcal{C}\ell}

\usepackage{wrapfig}
\usepackage{calc}%needed for the width/height calculations

\usepackage{pifont}

\usepackage{caption}
\usepackage{graphicx,lipsum}
\newcommand{\psf}{\ps^{\text{fin.}}}
\newcommand{\notimplies}{\;\not\!\!\!\implies}
\usepackage{mathdots}
\newcommand{\e}{\text{e}}
\usepackage{float}
\usepackage{wrapfig}
\DeclareMathOperator{\Th}{Th}
\DeclareMathOperator{\vl}{vl}
\DeclareMathOperator{\var}{var}
\DeclareMathOperator{\Var}{Var}
\DeclareMathOperator{\ar}{ar}
\newcommand{\Mydef}{\overset{\text{def}}{=}}
\DeclareMathOperator{\ED}{ED}
\DeclareMathOperator{\diag}{diag}
% Add this in the document preamble (main .tex)
\usepackage{bussproofs} % provides prooftree, \AxiomC, \UnaryInfC, \RightLabel


\DeclareMathOperator{\Ord}{Ord} % Order type
\DeclareMathOperator{\Dom}{Dom} % Domain of a function
\DeclareMathOperator{\Imm}{Im} % Image of a function
\DeclareMathOperator{\cof}{cof} % cofinality
\DeclareMathOperator{\Span}{Span} % Span of a set of vectors
\DeclareMathOperator{\supp}{supp} % support of a function
\DeclareMathOperator{\tc}{tc} % transitive closure
\newcommand{\A}{\mathbb{A}} % algebraic numbers
\newcommand{\psinf}{\mathcal{P}^{\text{inf.}}} % infinite power set
